
%%%%%%%%%%%%%%%%%%%%%%%%%%+ Lite review:%%%%%%%%%%%%%%%%%%%%%%%%%%%%%%%%%%%%%%%%%
%%%%% on attack generation (Leverage Ali's paper) (Ali）
Modern cyber-physical critical infrastructures (CPCI) are facing fast-evolving cyber threats. Cyberattacks can strongly impact the operation of the CPCIs. For instance, a suspected cyber intruder took control of the Prykarpattyaoblenergo power system control center in western Ukraine in December 2015, leaving 230,000 people without electricity for up to 6 hours. That was the first time hackers have successfully targeted a country's power grid \cite{liang20162015}. Also, cyber attacks in water systems have already become a reality. In 2015, 25 cyber attacks were disclosed in several water systems \cite{clark2017protecting}. And recently in 2020, a malicious cyber-attack attempted to raise the chlorine level in Israel's water supply to \textcolor{blue}{a} dangerous proportion \cite{brentan2021cyber}. Supervisory Control and Data Acquisition (SCADA) is a critical part of the CPCI that is highly susceptible to cyberattacks. \textcolor{blue}{It} is responsible for collecting measurements from Remote Terminal Units (RTU) and sending them to the Control Centers (CC) to perform various functions such as contingency analysis, optimal planning, state estimation (SE), etc. The primary purpose of those functions is to maintain stable and secure operation of the CPCI. 

SE is a core function in the CCs \textcolor{blue}{which} monitor\textcolor{blue}{s} the system states using information collected from the SCADA systems and other measurement devices. SE can also detect and identify bad data that can cause significant error in the resulting estimate. The SE algorithm's accuracy depends on how pure and accurate the system measurements are. If any malicious or erroneous measurements pass through undetected, it could mislead the operators into making catastrophically wrong decisions. \textcolor{blue}{Numerous studies have shown that traditional SE algorithms are vulnerable to False Data Injection Attacks} (\textcolor{blue}{FDIAs}), where an intruder aims to hack multiple RTUs or even communication channels to insert fake measurements \textcolor{blue}{in order to} misguide the CCs operational process \cite{liu2011false, boba, sou, kordestani2021observer}. Since \textcolor{blue}{FDIA} was introduced in \cite{liu2011false}, different researchers have proposed various techniques for detection and localization, including diagnostic robust generalized potential \cite{majumdar2016bad}, generalized likelihood ratio \cite{tang2016detection}, fast Go-Decomposition (GoDec) \cite{li2018detecting}, Markov chain \cite{li2017hmm}, Bayesian detection with binary hypothesis \cite{ding2018survey}, \textcolor{blue}{modified unbiased finite impulse response (UFIR) estimator \cite{ju2020deception}}, and cosine similarity matching scheme \cite{rawat2015detection} and some residual-based approaches such as the Kullback Leibler distance method \cite{chaojun2015detecting}, unscented Kalman filter \cite{vzivkovic2018detection}, $\chi^2$ failure detector \cite{mo2010false} and a residual-based localization scheme \cite{zheng2021Resilient}. However, \textcolor{blue}{due to their dependence} on the system model, the associated model uncertainty would affect \textcolor{blue}{their accuracy. Thus, if} the \textcolor{blue}{FDIA}s are \textcolor{blue}{designed appropriately,} these model-based or residual-based detectors can easily be bypassed \cite{liu2011false, zheng2021Resilient}. \textcolor{blue}{Take for example, filter-based attack detector designs \cite{ju2020deception, vzivkovic2018detection}; the recovered attack signals could have deviations due to the uncertainties of the underlying model used in the filter design. Thus, if the attacker substitutes current meter measurements with historical readings following the same distribution, the Bayesian framework and generalized likelihood ratio approaches are unable to identify this FDIA \cite{singh2017joint}.}

To overcome the limitations of traditional residual-based bad data detection approaches, data-driven solutions based on machine learning algorithms have been widely adopted for detecting and localizing \textcolor{blue}{FDIA}s due to their fast execution times and accurate results \cite{ozay2016machine, boyaci2021joint, esmalifalak2011stealth, svm2, wilson2018deep, yan2016detection, lee2001training}. \cite{ozay2016machine} utilized and tested various machine learning algorithms in detecting \textcolor{blue}{FDIA}s. The results showed that machine learning algorithms can detect \textcolor{blue}{FDIA}s accurately and faster. \cite{boyaci2021joint} proposed a deep neural network algorithm that can automatically detect and localize \textcolor{blue}{FDIA}s \textcolor{blue}{in power systems.} The proposed method can integrate the underlying graph topology of the grid and spatial correlations of its measurement data to jointly detect and localize the \textcolor{blue}{FDIA}s in power systems. \cite{lee2001training} proposed a hierarchy of neural networks with the capability of detecting novel attacks by training networks to understand the entire normal space instead of attack space.

%%%%%  Resilient estimators (Yu)
In addition to attack detection and localization, resilient estimation approaches focus on maintaining the correctness of state estimates in the presence of adversarial attacks, which enables resilient operation of the CPSs \cite{ding2020secure}. Although attacks are possibly unbounded and can be designed to "fool" $\ell_2$ observers \cite{mo2010false,liu2011false}, the attack signals are sparse due to the resource limitation of the attackers. Based on \textcolor{blue}{these} characteristics, \textcolor{blue}{the} authors in \cite{fawzi2014secure} proposed an $\ell_1$ minimization program for estimator design on linear systems and proposed an upper-bound on the total number of compromised measurements for guaranteed secure estimation. Some extension work considered robustness on resilient $\ell_1$ estimation in presence of modeling error \cite{pajic2014robustness}, worst-case estimation error bound in presence of bounded noise \cite{pajic2015attack}, and $\ell_1-\ell_2$ estimation scheme \cite{lee2015secure}. And some iterative resilient observer designs has been proposed for determinant system, such as event-trigger Luenburger Observer \cite{shoukry2015event} and Gramian-based estimator \cite{chong2015observability}. For stochastic system, several attack-resilient control \textcolor{blue}{strategies} have also been studied \cite{ding2016security, wan2019attack}. \textcolor{blue}{In addition, resilient estimation and control design has also been studied for nonlinear systems \cite{wang2018resilient}.  }
% using Takagi-Sugeno (T-S) fuzzy modeling method \cite{xie2017relaxed}

%%%%%% Improved resilient estimation (Yu)
\textcolor{blue}{Furthermore, some enhanced resilient CPS designs have been studied to accommodate even worse adversarial condition. To push the} limit on \textcolor{blue}{the} number of compromised sensors below which secure estimation would be guaranteed, \textcolor{blue}{additional prior information could be used} in the resilient estimation design. Existing results have used such information as measurement prior, support prior, and state prior\textcolor{blue}{: In \cite{shinohara2019resilient}, three forms of \emph{state prior} are considered: \emph{sparsity information}, \emph{$(\alpha, \overline{n}_0)$ sparsity information,}, and \emph{side information.} These are information generally unavailable to the malicious agent. In \cite{wan2019attack}, the authors designed an improved resilient estimator, with input and state constraints, against actuator attacks.} In our previous work \cite{anubi2019enhanced, anubi2019resilient,anubi2020multi}, \textcolor{blue}{the} resiliency of \textcolor{blue}{$\ell_1$ based} observers is enhanced by utilizing a measurement prior, generated by \textcolor{blue}{a pre-trained Gaussian process}. \textcolor{blue}{In} \cite{zheng2020attack}, a resilient unscented Kalman filter was proposed based on improved support prior. In \cite{zheng2021Resilient}, a weighted $\ell_1$ observer is designed to exploit the support prior more rigorously.

%%%%%% Automated attack generation (Yu)
 \textcolor{blue}{Alternatively, improved resiliency of CPS could also be achieved by using more accurate attack detection and localization algorithms. For data-driven attack detectors, especially supervised learning-based detectors, exploiting more of the attack space and system nominal conditions could improve accuracy significantly \cite{sayghe2020survey}. This has motivated researchers to study attack data generation strategies for improving the detection/localization boundaries of the underlying data-driven algorithms. The goal of the attack generation problem is to generate feasible attacks that could elicit significant targeted disruption on system, while bypassing bad data detection. Model-based attack generation strategies often assume the attackers have the knowledge of system models \cite{liu2011false,mo2010false}. This assumption limits the application range of those attack generation strategies. Thus, automated attack generation without knowledge of the system models, is currently gaining a lot of attention. Although this topic has been studied widely in the machine learning area, for example; adversarial machine learning \cite{huang2011adversarial} and deepfake \cite{westerlund2019emergence}, the automated attack generation problem for CPS is a new challenging problem since the physical process dynamics is included. Inspired by the generative adversarial network (GAN) \cite{goodfellow2014generative}, some earlier researches \cite{mohammadpourfard2020generation, ahmadian2018cyber, shahriar2021iattackgen} have studied GAN-based attack generation by using an existing attack dataset. The idea is to use GAN to learn the underlying distribution of feasible attacks, then a richer dataset of attacks could be generated. However, it is not practical to assume the existence of a non-generative attack dataset. Thus, an automated attack generation process, without attack examples, is needed. To the best of the authors knowledge, this paper is the first paper to study this problem.}
 
 % Other security assessment researches could also be found in \cite{kundur2011towards, srivastava2013modeling, kwon2013security, ten2008vulnerability}.


%%%%%%%%%%%%%%%%%%%%%%%%%%% Novelty argument %%%%%%%%%%%%%%%%%%%%%%%%%%%%%%%%%%%%%%%% (Yu)

%%%%% What is new?
% a.GAN method to generate true \textcolor{blue}{FDIA}  
% b. Pruning to improve resiliency  
% c. Resilient observer based control design
%%%%% Improvement over existing methods: The advantages of the approach over all lit. reviewed papers above.
\emph{Contribution}: This paper proposes an enhanced resilient solution for CPSs subject to FDIAs. It is a hybrid physics-based and data-driven, detection-based resilient estimation framework built around a \emph{Kalman filter} with a \emph{chi-square} innovation detector. The framework consists of an automated attack generator, supervised learning-based attack detection and localization, support pruning, and a resilient KF on the retained safe sensors. To train the supervised detector, abundant adversarial examples are required. Rather than relying on expensive optimization-based generators \cite{zheng2021Resilient, yang2013false} or labeled attack corpora \cite{mohammadpourfard2020generation,ahmadian2018cyber}, we propose an unsupervised generative model trained by physics-based discriminators derived from the KF innovation (stealthiness against $\chi^2$) and the closed-loop critical-output map (effectiveness). Secondly, a pruning algorithm systematically improves localization precision, with guarantees quantified as a function of the pruning reliability level. Thirdly, we establish a \emph{quantitative bridge} among attack sparsity, post-pruning observability, KF mean-square estimation error (via the subset Riccati covariance), and detection/localization precision ($\textsf{PPV}_{\eta}$ and the $\chi^2$ false-alarm specification). Consequently, the CPS can maintain certified estimation performance while attacks are underway.


%%%%%%%%%%%%%%%%%%%%%%%%%%%% Paper organization and other stuff %%%%%%%%%%%%%%%%%%%%%%%%%%%%%%% (Yu)
The remainder of the paper is organized as follows. Section~\ref{sec:notation} collects notation. Section~\ref{sec:system_model} introduces the stochastic CPS model, the Kalman filter, and the chi-square detector. Section~\ref{sec:attack_generation} defines successful FDIAs against this pair and presents the physics-guided generative model. Section~\ref{sec:detection_localization} develops MLP-based detection/localization, pruning, and the sparsity--observability--precision bridge. Section~\ref{sec:resilient_estimator} designs the resilient KF and links subset covariances to pruning reliability. Section~\ref{sec:simulation} reports water-distribution simulations, and Section~\ref{sec:conclusion} concludes.}